% Generated draft: requires mathematical and source review.
\documentclass[11pt]{article}
\usepackage[margin=25mm]{geometry}
\usepackage{fontspec}
\setmainfont{Latin Modern Roman}
\usepackage{amsmath,amssymb,mathtools,unicode-math}
\setmathfont{Latin Modern Math}
\usepackage{enumitem,needspace,xurl,hyperref}
\setlength{\parindent}{0pt}
\setlength{\parskip}{5pt}
\setlength{\emergencystretch}{3em}
\newcommand{\N}{\mathbb{N}}
\newcommand{\Z}{\mathbb{Z}}
\newcommand{\Q}{\mathbb{Q}}
\newcommand{\R}{\mathbb{R}}
\newcommand{\C}{\mathbb{C}}
\newcommand{\F}{\mathbb{F}}
\newcommand{\A}{\mathbb{A}}
\newcommand{\PP}{\mathbb{P}}
\newcommand{\E}{\mathbb{E}}
\newcommand{\eps}{\varepsilon}
\newcommand{\dd}{\,\mathrm{d}}
\begin{document}
\section*{UnsolvedMath ID: 30002971}
\textbf{Source problem code:} catalogue entry
\subsection{Definitions and mathematical statement}
\textbf{Scope and terminology.} The source concerns morphogenesis in a tissue whose cells regulate their behavior locally, while tissue growth or patterning may require information about overall tissue size. Here, ``local regulation'' and ``global tissue-size information'' are descriptive source terms: the record specifies neither the information available to each cell nor a mathematical measure of tissue size. Active mechanical stress refers to mechanically generated signaling associated with cellular activity; its constitutive definition is not supplied. The expression ``its derivative'' does not specify a spatial derivative, a time derivative, or another operation.

\textbf{Source question.} How can information about global tissue size be transmitted to locally regulated cells? In particular, can active mechanical stress, or an appropriate derivative of that stress, provide a locally accessible signal that depends on overall tissue size and conveys size information beyond the fixed characteristic length attributed in the source to reaction--diffusion mechanisms?

\textbf{Unresolved mathematical specification.} This is a qualitative modeling question, not a fully specified conjecture. A precise mathematical formulation requires a tissue domain and its dimension, a definition of size, mechanical and biochemical evolution laws, admissible local observations, initial and boundary conditions, and a criterion for successful transmission or recovery of size information. The record supplies none of these. It also does not specify whether the target is existence of one mechanism, identification of a biologically relevant mechanism, or a statement valid for a class of models, nor whether transmission is transient or stationary. Consequently, no domain, quantifier over models, or boundary condition is imposed here. The comparison with a fixed reaction--diffusion characteristic length is retained as a source-specific qualification; it is not asserted as a theorem about all reaction--diffusion systems.
\subsection{Short English statement}
The question asks how cells acting on local information can learn the overall size of a developing tissue, particularly through active mechanical stress or a derivative of it. The source does not supply enough mathematical structure to turn this into a uniquely quantified conjecture. Its supplied open status requires independent verification.
\subsection{Sources}
\begin{itemize}
\item Supplied local UnsolvedMath registry record, permanent ID \(30002971\), collection \emph{Oberwolfach Reports Open Problems}, exact source code \texttt{catalogue entry}, title \emph{Long-Range Mechanical Signaling in Tissue Morphogenesis}. No individual proposer is named. The record gives proposed year 2015.
\item Supplied report attribution: \emph{Mini-Workshop: Mathematics of Differential Growth, Morphogenesis, and Pattern Selection} (2016), Oberwolfach Reports, pp.\ 2895--2910, DOI \url{https://doi.org/10.4171/owr/2015/49}. The supplied background describes Oberwolfach Reports as documenting mathematical workshops and their open-problem discussions. The relationship between the proposed year 2015, the citation year 2016, and the DOI containing 2015 has not been independently checked; no problem-specific page or speaker attribution is supplied.
\item Supplied literature reference: Ambrosi--Pettinati--Ciarletta, \emph{Active Stress}, \url{https://doi.org/10.1016/j.ijnonlinmec.2014.11.027}.
\item Supplied literature reference: Collinet--Lecuit, \emph{Programmed and Self-Organized Information Flow}, \url{https://doi.org/10.1038/nrm.2017.139}.
\item Supplied literature reference: \emph{Measuring Multiscale Mechanical Information}, \url{https://doi.org/10.1016/j.gde.2011.08.008}. No authors are given for this reference in the record.
\item Literature assessment dated exactly \texttt{2026-08-21}: the supplied curation labels the problem open and describes mechanochemical stress, morphogen gradients, and biochemical feedback as plausible transmission mechanisms, while stating that the general question of how local cells acquire global size information remains unresolved. This is supplied dated curation derived from the OWR record, not an independently verified literature assessment. The background lists the three references above as sources checked by that curation; no source visit or verification is claimed here.
\end{itemize}
\end{document}
